%%%%%%%%%%%%%%%%%%%%%%%%%%%%%%%%%%%%%%%%%%%%%%%%%%%%%%%%%%%%%
% L. D. Abreu, \'O. Ciaurri and J. L. Varona,               %
% The spectrum of the right inverse of the Dunkl operator   %
%                                                           %
% LAST MODIFICATION DATE: 10 - November- 2012               %
%                                                           %
% Written with LaTeX, class amsart                          %
%%%%%%%%%%%%%%%%%%%%%%%%%%%%%%%%%%%%%%%%%%%%%%%%%%%%%%%%%%%%%


%%%%%%%%%%%%%%%%%%%%%%%%%%%%%%%%%%%%%%%%%%%%%%%%%%%%
\documentclass[12pt,a4paper]{amsart}
%%%%%%%%%%%%%%%%%%%%%%%%%%%%%%%%%%%%%%%%%%%%%%%%%%%%

%--------------------------------------------------------
%\usepackage[notcite,notref]{showkeys} % <----- For drafts
%\usepackage{color}
%--------------------------------------------------------

\usepackage{amsfonts}

\usepackage{url}


%% FOR DRAFTS AT 12pt:
\advance\textwidth by 1.5cm

% \usepackage{backref} % for drafts

\theoremstyle{plain}
\newtheorem{Thm}{Theorem}
%\newtheorem{Cor}[Thm]{Corollary}
%\newtheorem*{Prop}{Proposition}
%\newtheorem{Lem}{Lemma}%[section]

\theoremstyle{definition}
%\newtheorem{Defn}{Definition}
%\newtheorem*{Ackn}{Acknowledgement}
%\theoremstyle{remark}
\newtheorem*{Rem}{Remark}

%\newcommand{\supp}{\operatorname{supp}}
%\newcommand{\Id}{\operatorname{Id}}
\newcommand{\spa}{\operatorname{span}}
%\renewcommand{\Im}{\operatorname{Im}}
%\renewcommand{\Re}{\operatorname{Re}}

\newcommand{\N}{\mathbb{N}}
\newcommand{\Z}{\mathbb{Z}}
\newcommand{\R}{\mathbb{R}}
\newcommand{\C}{\mathbb{C}}
%\newcommand{\calK}{\mathcal{K}}
%\newcommand{\calS}{\mathcal{S}}
\newcommand{\calF}{\mathcal{F}}
\newcommand{\calI}{\mathcal{I}}
%\newcommand{\calJ}{\mathcal{J}}
%\newcommand{\calP}{\mathcal{P}}
%\newcommand{\calQ}{\mathcal{Q}}


\allowdisplaybreaks[3]


%------------------------------
\begin{document}
%------------------------------

%\date{November 9, 2012.}
\renewcommand{\datename}{\textbf{THIS PAPER HAS BEEN PUBLISHED IN:}}
\date{\textit{Rev.\@ Mat.\@ Complut.\@} \textbf{26} (2013), 471--483.}

\title[The spectrum of the right inverse of the Dunkl operator]{%
The spectrum of the right inverse of the Dunkl operator}


\author[L. D. Abreu]{Lu\'{\i}s Daniel Abreu}
\address{CMUC and Departamento de Matem\'atica,
Universidade de Coimbra,
Faculdade de Ci\^{e}ncias e Tecnologia (FCTUC),
3001-454 Coimbra, Portugal}
\curraddr{Austrian Academy of Sciences, Acoustic Research Institute, 
Reichsratsstrasse 17, A-1010 Vienna, Austria}
\email{daniel@mat.uc.pt}
\thanks{Research of the first author supported by FCT project
SPTDC/MAT/114394/2009, POCI 2010 and FSE, and Austrian Science Fund (FWF) 
START-project FLAME (``Frames and Linear Operators for Acoustical Modeling
and Parameter Estimation'', Y 551-N13)}

\author[\'O. Ciaurri]{\'Oscar Ciaurri}
\address{CIME and Departamento de Matem\'aticas y Computaci\'on,
Universidad de La Rioja, % Edificio J.~L.~Vives,
% Calle Luis de Ulloa s/n,
26004 Logro\~no, Spain}
\email{oscar.ciaurri@unirioja.es}

\author[J. L. Varona]{Juan Luis Varona}
\address{CIME and Departamento de Matem\'aticas y Computaci\'on,
Universidad de La Rioja, % Edificio J.~L.~Vives,
% Calle Luis de Ulloa s/n,
26004 Logro\~no, Spain}
\email{jvarona@unirioja.es}
\urladdr{http://www.unirioja.es/cu/jvarona/}
\thanks{Research of the second and third authors supported 
by grant MTM2012-36732-C03-02 of the DGI}


\keywords{Plane wave expansion, Dunkl operator, Dunkl kernel, spectrum, 
diagonalization of operators, special functions.}
\subjclass[2010]{Primary 42C10; Secondary 47A10, 47A75, 33C10, 33C45}


\begin{abstract}
From the Dunkl analogue of Gegenbauer's expansion of the plane wave, 
we derive an explicit closed formula for the spectrum of a right 
inverse of the Dunkl operator. This is done by stating the problem 
in such a way it is possible to use the technique due to Ismail and Zhang.
\end{abstract}

%------------------------------
\maketitle
%------------------------------

%\tableofcontents

%------------------------------
\section{Introduction}
%------------------------------

During the last years, and starting in the seminal paper~\cite{IZ}, 
Mourad Ismail and his collaborators had studied the diagonalization 
of the inverse of several differential or difference operators. 
In particular, it has been done in \cite{IZ, IRZ} for the differential 
and the Askey-Wilson divided difference operators in different polynomial 
bases, in~\cite{IR} for the full Askey-Wilson weight,
in~\cite{ISem} for the $q$-difference operator, and in~\cite{IJacob} 
for the $q^{-1}$-Askey-Wilson operator.

This article is devoted to the diagonalization of the right inverse of 
the Dunkl operator on the real line given by
\[
  \Lambda_\alpha f(x) = \frac{d}{dx}f(x) + \frac{2\alpha+1}{x}
  \left(\frac{f(x)-f(-x)}{2}\right).
\]
The operator $\Lambda_\alpha$ is a generalization of the derivative 
$\frac{d}{dx}$ (which is the case corresponding to $\alpha=-1/2$) 
whose study, as well as of the so-called Dunkl transform that is 
closely related to~$\Lambda_\alpha$, has engender many recent papers 
(see, for instance,~\cite{AJ, BCV, CV, Dunkl, Dunkl2, Jeu, Rosenblum, Rosler}). 
In the present paper, we are going to find the spectrum of a right inverse 
for the Dunkl operator $\Lambda_\alpha$ in spaces weighted by the appropriated 
weight functions.

In~\cite{IZ}, an important tool to study the inverse of the differential 
operator $\frac{d}{dx}$ is Gegenbauer's expansion of the plane wave in 
ultraspherical polynomials and Bessel functions 
(see~\cite[\S\,4.8, formula (4.8.3), p.~116]{Ism}):
\begin{equation}
\label{eq:geg}
  e^{ixt} = \Gamma(\beta) \left(\frac{x}{2}\right)^{-\beta}
  \sum_{n = 0}^{\infty} i^{n}(\beta+n)J_{\beta+n}(x)C_{n}^{\beta}(t)
\end{equation}
(in the particular case $\beta=0$, this formula is the so-called
Jacobi-Anger identity). Instead of this expansion for $e^{ixt}$,
that is the kernel of the Fourier transform, we will require an
expansion for $E_{\alpha}(ixt)$, the kernel of the Dunkl transform
on the real line. Thus, we will use that the Dunkl kernel
$E_{\alpha}(ixt)$ can be expanded as
\begin{equation}
\label{eq:Dunklexpans}
  E_\alpha(ixt) = \Gamma(\alpha+\beta+1) \left(\frac{x}{2}\right)^{-\alpha-\beta-1}
  \sum_{n=0}^\infty i^n (\alpha+\beta+n+1)J_{\alpha+\beta+n+1}(x)
  C_n^{(\beta+1/2,\alpha+1/2)}(t),
\end{equation}
where $C_n^{(\beta+1/2,\alpha+1/2)}$ are the so-called generalized
Gegenbauer polynomials and $J_a$ denotes the Bessel function of
order~$a$. The identity~\eqref{eq:Dunklexpans} holds for $\alpha,\beta>-1$ 
and $\alpha+\beta>-1$. It was proved by M.~R\"{o}sler in~\cite{Rosler} 
for $\alpha\ge -1/2$, and by the authors in~\cite{ACV} for $\alpha>-1$. 
By using~\eqref{eq:Dunklexpans} in the place of~\eqref{eq:geg}, a
part of the technique for $\frac{d}{dx}$ can be followed, and the
results obtained are very satisfactory, finally getting a nice
extension of the result of Ismail and Zhang in~\cite{IZ}, in the
same sense as the Bessel functions generalize the trigonometric
functions, the Dunkl operator has the derivative as a particular
case, or the Dunkl transform is an extension of the Fourier
transform.

The structure of the paper is as follows.
In Section~\ref{sec:prelim} we state the precise notation and give some 
preliminaries related to the Dunkl operator on the real line and to the 
generalized Gegenbauer polynomials, as well as the identity~\eqref{eq:Dunklexpans}, 
that is the main tool to solve our problem.
The essential part of the paper is included in Section~\ref{sec:inverse}.
We define the right inverse of the Dunkl operator in spaces weighted by 
the weight function of the generalized Gegenbauer polynomials and use 
this definition to find explicit formulas for the spectrum of this 
right inverse that we will denote~$T_{\beta,\alpha}$. It turns out that 
the eigenvalues of the resulting integral operator $T_{\beta,\alpha}$ 
are defined in terms of zeros of Bessel functions 
$\{\pm i/j_{\alpha+\beta+1,k}\}_{k\ge 1}$ and the eigenfunctions are
\[
  g_{\pm i/j_{\alpha+\beta+1,k}}(t) = \mp i
  \left(\frac{j_{\alpha+\beta+1,k}}{2}\right)^{\alpha+\beta+1}
  \frac{E_\alpha(\mp t j_{\alpha+\beta+1,k})}
    {\Gamma(\alpha+\beta+1)(\alpha+\beta+2)J_{\alpha+\beta}(j_{\alpha+\beta+1,k})};
\]
actually, this is the main result of the paper, whose complete details can be seen 
in Theorem~\ref{thm:main}.

%------------------------------
\section{Notation and preliminaries}
\label{sec:prelim}
%------------------------------

%------------------------------
\subsection{The Dunkl operator on the real line}
%------------------------------

For $\alpha > -1$, let $J_\alpha$ denote the Bessel function of order $\alpha$ and,
for complex values of the variable $z$, let
\begin{equation*}
  \calI_\alpha(z)
  = 2^\alpha \Gamma(\alpha+1)\, \frac{J_\alpha(iz)}{(iz)^\alpha}
  = \Gamma(\alpha+1) \sum_{k=0}^\infty \frac{(z/2)^{2k}}{k!\,\Gamma(k+\alpha+1)}
\end{equation*}
($\calI_\alpha$ is a small variation of the so-called modified
Bessel function of the first kind and order~$\alpha$, usually
denoted by $I_\alpha$; see~\cite{Wat}). Moreover, let us take
\begin{equation*}
  E_\alpha(z) = \calI_\alpha(z) + \frac{z}{2(\alpha+1)}\,
    \calI_{\alpha+1}(z),
  \qquad z \in \C;
\end{equation*}
this function is called the Dunkl kernel.

The Dunkl operators on $\R^n$ are differential-difference operators
associated with some finite reflection groups (see~\cite{Dunkl}). We
consider the Dunkl operator $\Lambda_\alpha$, $\alpha \ge -1/2$, associated
with the reflection group $\Z_2$ on $\R$ given by
\begin{equation}
\label{eq:dunkloper}
  \Lambda_\alpha f(x) = \frac{d}{dx}f(x) + \frac{2\alpha+1}{x}
  \left(\frac{f(x)-f(-x)}{2}\right).
\end{equation}
For $\alpha \ge -1/2$ and $\lambda\in \C$, the initial value problem
\begin{equation}
\label{eq:dunklproblem}
\begin{cases}
  \Lambda_\alpha f(x) = \lambda f(x),\quad x\in \R, \\
  f(0) = 1
\end{cases}
\end{equation}
has $E_\alpha(\lambda x)$ as its unique solution (see~\cite{Dunkl2}
and~\cite{Jeu}). For $\alpha = -1/2$, it is
clear that $\Lambda_{-1/2} = d/dx$, and $E_{-1/2}(\lambda x) = e^{\lambda x}$.

Along this paper we will usually employ $E_\alpha$ in the form
\begin{equation}
\label{eq:Ea}
  E_{\alpha}(ix) = 2^{\alpha}\Gamma(\alpha+1)
  \left( \frac{J_{\alpha}(x)}{x^{\alpha}}+\frac{J_{\alpha+1}(x)}
    {x^{\alpha+1}}\,xi\right).
\end{equation}
If, moreover, we take
\[
  d\mu_{\alpha}(x) = (2^{\alpha+1} \Gamma(\alpha+1))^{-1}|x|^{2\alpha+1}\,dx,
\]
the name Dunkl kernel for $E_\alpha$ is justified because it is the kernel 
of the Dunkl transform of order $\alpha \geq -1/2$ given by
\begin{equation}
\label{eq:Du-T}
  \calF_{\alpha}f(x) = \int_{\R} f(t)E_{\alpha}(-ixt)
  \,d\mu_{\alpha}(t),\quad x\in \R,
\end{equation}
for $f\in L^{1}(\R,d\mu_{\alpha})$.
In a similar way to the Fourier transform (which is the particular
case $\alpha = -1/2$), the definition is extended $L^{2}(\R,d\mu_{\alpha})$ 
and $\calF_{\alpha}$ becomes to be an isometric isomorphism on $L^{2}(\R,d\mu_{\alpha})$.
Yet more, the Dunkl transform $\calF_{\alpha}$ can also be defined in
$L^2(\R,d\mu_{\alpha})$ for $-1 < \alpha < -1/2$; although in this case the
expression~\eqref{eq:Du-T} is no longer valid for
$f\in L^1(\R,d\mu_{\alpha})$ in general, it preserves the same properties in
$L^2(\R,d\mu_{\alpha})$ (see~\cite{Rosenblum} for details).
Following this spirit, we will do our study not only for $\alpha\ge-1/2$, 
but for the whole range $\alpha > -1$.

%------------------------------
\subsection{Generalized Gegenbauer polynomials}
%------------------------------

Following~\cite[Definition~1.5.5, p.~27]{DX}, let us introduce the generalized
Gegenbauer polynomials $C^{(\lambda,\nu)}(t)$ for $\lambda>-1/2$, $\nu \ge 0$
and $n \ge 0$
(the case $\nu=0$ corresponding with the ordinary Gegenbauer polynomials);
actually, for convenience with the notation of this paper,
we are going to use $C_n^{(\beta+1/2,\alpha+1/2)}(x)$. In this way, for $\beta>-1$
and $\alpha\ge -1/2$, the generalized Gegenbauer polynomials are defined by
\begin{align}
\label{eq:C2n}
  C_{2n}^{(\beta+1/2,\alpha+1/2)}(t)
  &= (-1)^n \, \frac{(\alpha+\beta+1)_n}{(\alpha+1)_n} P_n^{(\alpha,\beta)}(1-2t^2), \\
\label{eq:C2n1}
  C_{2n+1}^{(\beta+1/2,\alpha+1/2)}(t)
  &= (-1)^n \, \frac{(\alpha+\beta+1)_{n+1}}{(\alpha+1)_{n+1}} t P_n^{(\alpha+1,\beta)}(1-2t^2),
\end{align}
where in the coefficients we are using the Pochhammer symbol
$(a)_n = a(a+1)\cdots(a+n-1) = \Gamma(a+n)/\Gamma(a)$. 
Note that there is no problem in extending the definition of the generalized Gegenbauer 
polynomials taking $\alpha>-1$, so we will assume this situation.

The definitions~\eqref{eq:C2n} and~\eqref{eq:C2n1} easily allow to translate the 
well known orthogonality of the Jacobi polynomials on $(-1,1)$ to the orthogonality 
of the generalized Gegenbauer polynomials on $(-1,1)$ with a related weight, 
which is more useful in the context of this paper.
In particular, the generalized Gegenbauer polynomials 
$\{C_{k}^{(\beta+1/2,\alpha+1/2)}(t)\}_{k\ge0}$ form a complete 
orthogonal system on $L^2((-1,1),(1-t^2)^\beta \,d\mu_\alpha(t))$ 
for $\alpha,\beta>-1$.
Moreover, from the $L^2$-norm of the Jacobi polynomials
(see~\cite[formula 22.2.1, p.~774]{AbSt}),
it is easy to find
\begin{align}
\label{eq:hC2n}
  h_{2n}^{(\beta,\alpha)}
  &= \int_{-1}^1 \left[C_{2n}^{(\beta+1/2,\alpha+1/2)}(t)\right]^2
  (1-t^2)^\beta \,d\mu_\alpha(t) \\
  &= \frac{1}{2^{\alpha+1}} \frac{\Gamma(\alpha+1)\Gamma(\beta+n+1)\Gamma(\alpha+\beta+n+1)}
    {(\alpha+\beta+2n+1)\Gamma(\alpha+\beta+1)^2 \Gamma(\alpha+n+1)\, n!},
  \notag \\
\label{eq:hC2n1}
  h_{2n+1}^{(\beta,\alpha)}
  &= \int_{-1}^1 \left[C_{2n+1}^{(\beta+1/2,\alpha+1/2)}(t)\right]^2
  (1-t^2)^\beta \,d\mu_\alpha(t) \\
  &= \frac{1}{2^{\alpha+1}} \frac{\Gamma(\alpha+1)\Gamma(\beta+n+1)\Gamma(\alpha+\beta+n+2)}
    {(\alpha+\beta+2n+2) \Gamma(\alpha+\beta+1)^2 \Gamma(\alpha+n+2)\, n!}.
\notag
\end{align}

%------------------------------
\subsection{The expansion of the Dunkl kernel}
%------------------------------

The last ingredient in this section with preliminaries is to explain 
the~\eqref{eq:Dunklexpans} that we give in the introduction. 
Firstly, let us note that, when you use something like ${J_{a+n}(x)}/{x^a}$, we mean
\[
  \frac{J_{a+n}(x)}{x^a}
  = \sum_{k=0}^\infty \frac{(-1)^k}{2^ k!\,\Gamma(k+a+n+1)}
  \Bigl(\frac{x}{2}\Bigr)^{2k+n},
\]
that is an analytic function on the complex plane for every $a>-1$ 
and every integer $n \ge 0$.

Then, the precise statement of the Dunkl analogue of Gegenbauer's 
expansion of the plane wave is that, for $\alpha,\beta > -1$ and $\alpha+\beta>-1$, 
the Dunkl kernel can be expanded, for $x\in \R$, as
\begin{equation}
\label{eq:JA-Dunkl}
  E_\alpha(ixt) = 2^{\alpha+\beta+1} \Gamma(\alpha+\beta+1)
  \sum_{n = 0}^{\infty} i^{n}(\alpha+\beta+n+1)
  \frac{J_{\alpha+\beta+n+1}(x)}{x^{\alpha+\beta+1}} C_{n}^{(\beta+1/2,\alpha+1/2)}(t),
\end{equation}
and this expansion holds in $L^{2}((-1,1),d\mu_{\alpha}(t))$.

This important identity has been proved in~\cite{Rosler} for $\alpha\ge -1/2$ 
by showing that, under this restriction, formulas~\eqref{eq:geg} and~\eqref{eq:JA-Dunkl} 
are equivalent. The proof in one direction is clear, just by specializing the parameters. 
To obtain~\eqref{eq:JA-Dunkl} from~\eqref{eq:geg}, we can use the intertwining operator
\begin{equation*}
  V_\alpha g(t) = \frac{\Gamma(2\alpha+2)}{2^{2\alpha+1} \Gamma(\alpha+1/2)\Gamma(\alpha+3/2)}
  \int_{-1}^1 g(st) (1-s)^{\alpha-1/2} (1+s)^{\alpha+1/2}\,ds
\end{equation*}
(see~\cite[Definition~1.5.1, p.~24]{DX}, we change the parameter $\mu$ 
in the definition given in~\cite{DX} by $\alpha+1/2$), defined for $\alpha\ge -1/2$. 
With this notation we have
\begin{equation*}
  V_\alpha C_n^{\alpha+\beta+1}(t) = C_n^{(\beta+1/2,\alpha+1/2)}(t)
\end{equation*}
and
\begin{equation*}
  V_\alpha e^{i\cdot}(t) = E_\alpha(it).
\end{equation*}
In this way, applying $V_\alpha$ to~\eqref{eq:geg} 
(with $\alpha+\beta+1$ instead of~$\beta$) we get~\eqref{eq:JA-Dunkl}.

This idea cannot be used in the whole range $\alpha>-1$, 
and the extension of~\eqref{eq:JA-Dunkl} to this case can be found in~\cite{ACV}.

%\begin{Rem}
%Another way of obtaining~\eqref{eq:JA-Dunkl} is as follows. Some results from~\cite{FW}
%were generalizated in~\cite{Ve} to
%\begin{multline*}
%  \sum_{m=0}^{\infty} a_mb_m \frac{(zw)^m}{m!}
%  = \sum_{n=0}^{\infty} \frac{(-z)^n}{n!\, (\gamma+n)_n}
%  \left( \sum_{r=0}^\infty \frac{b_{n+r}z^r}{r!\,(\gamma+2n+1)_r} \right) \\
%  \times \left( \sum_{s=0}^n \frac{(-n)_s (n+\gamma)_s}{s!} \,a_s w^s \right).
%\end{multline*}
%When $z$ and $w$ are replaced by $z\gamma$ and $w/\gamma$, respectively, 
%and we let $\gamma \to \infty$, %we get the companion formula
%\begin{equation*}
%  \sum_{m=0}^{\infty} a_mb_m (zw)^m
%  = \sum_{n=0}^{\infty} \frac{(-z)^n}{n!}
%  \Biggl( \sum_{j=0}^\infty \frac{b_{n+j}}{j!} \, z^j \Biggr)
%  \Biggl( \sum_{k=0}^n (-n)_k a_k w^k \Biggr)
%\end{equation*}
%(these formulas are also stated in~\cite[Chapter 9]{Ism}).
%These expansions contain several expansions in terms of Jacobi 
%polynomials of argument $1 - 2t^2$ (i.e., generalized Gegenbauer polynomials). 
%In particular, \eqref{eq:JA-Dunkl} follows in this way.
%\end{Rem}

%------------------------------
\section{Diagonalization of the right inverse of Dunkl operator}
\label{sec:inverse}
%------------------------------

As stated in the introduction, Ismail and Zhang study in~\cite{IZ} the
eigenfunctions and the eigenvalues of the right inverse of the derivative operator,
and the main tool is a suitable expansion of the corresponding plane wave.
The aim of this section is the analysis of this question for the
Dunkl operator~\eqref{eq:dunkloper} using the expansion~\eqref{eq:JA-Dunkl}.
This is a very natural extension, because the ordinary derivative
corresponds to the Dunkl operator $\Lambda_\alpha$ with $\alpha=-1/2$.

To simplify the notation, let us introduce the measure
\begin{equation*}
  d\mu_{\beta,\alpha}(t) = (1-t^2)^\beta \, d\mu_\alpha(t)
  = \tfrac{1}{2^{\alpha+1} \Gamma(\alpha+1)} (1-t^2)^\beta |t|^{2\alpha+1}\,dt.
\end{equation*}

It is not difficult to prove, using the appropriate recurrence
relations for the Jacobi polynomials, that the Dunkl operator over
the generalized Gegenbauer operator satisfies
\begin{equation}
\label{eq:LCn}
  \Lambda_\alpha C_n^{(\beta+1/2,\alpha+1/2)} (t)
  = 2(\alpha+\beta+1) C_{n-1}^{(\beta+3/2,\alpha+1/2)} (t).
\end{equation}

Motivated by the identity~\eqref{eq:LCn}, we define the right
inverse operator of $\Lambda_\alpha$ over $L^2((-1,1),
d\mu_{\beta+1,\alpha})$ by
\begin{equation*}
  T_{\beta,\alpha} g(t) = \sum_{n=1}^{\infty}
  \frac{g_{n-1}}{2(\alpha+\beta+1)} \, C_n^{(\beta+1/2,\alpha+1/2)}(t),
\end{equation*}
whenever $g$ has the expansion
\begin{equation*}
  g(t) = \sum_{n=0}^{\infty} g_n C_n^{(\beta+3/2,\alpha+1/2)} (t),
\end{equation*}
with
\begin{equation*}
  g_n = \big( h_n^{(\beta+1,\alpha)} \big)^{-1} \int_{-1}^1 g(r)
  C_n^{(\beta+3/2,\alpha+1/2)} (r) \,d\mu_{\beta+1,\alpha}(r)
\end{equation*}
(the norms $h_n^{(\beta+1,\alpha)}$ are given in~\eqref{eq:hC2n} and~\eqref{eq:hC2n1}).
Then, we can write
\begin{equation*}
  T_{\beta,\alpha} g(t) = \int_{-1}^1 g(r) K_{\beta,\alpha}(t,r) \,
  d\mu_{\beta+1,\alpha}(r),
\end{equation*}
with
\begin{equation*}
  K_{\beta,\alpha}(t,r)
  = \frac{1}{2(\alpha+\beta+1)}\sum_{n=1}^{\infty}
  \frac{C_n^{(\beta+1/2,\alpha+1/2)}(t) C_{n-1}^{(\beta+3/2,\alpha+1/2)}(r)}
    {h_{n-1}^{(\beta+1,\alpha)}}.
\end{equation*}

We want to diagonalize $T_{\beta,\alpha}$. To this end we have to
find values $\lambda \in \C$ and functions $g_{\lambda} \in
L^2((-1,1), d\mu_{\beta,\alpha}) \cap L^2((-1,1),
d\mu_{\beta+1,\alpha})$ such that
\begin{equation}
\label{eq:problem-auto}
  \lambda g_\lambda(t) = T_{\beta,\alpha} g_\lambda(t),
\end{equation}
where the expansion of $g$ in
$\{C_n^{(\beta+1/2,\alpha+1/2)}\}_{n\ge 0}$ is
\begin{equation}
\label{eq:expaneigen}
  g_\lambda(t) \sim \sum_{n=1}^{\infty}
  a_n(\lambda) C_n^{(\beta+1/2,\alpha+1/2)}(t).
\end{equation}


To find the eigenvalues and the eigenfunctions we start obtaining
a recurrence relation for the coefficients $\{a_n(\lambda)\}_{n\ge 1}$.
This is done by using an expression to write
$(1-r^2)C_{n-1}^{(\beta+3/2,\alpha+1/2)}(r)$ in terms of
$C_{n-1}^{(\beta+1/2,\alpha+1/2)} (r)$ and
$C_{n+1}^{(\beta+1/2,\alpha+1/2)} (r)$. This relation can be
deduced from the identity for the Jacobi polynomials
(see \cite[formula 22.7.16, p.~782]{AbSt})
\begin{equation}
\label{eq:Jacobirecu}
  P_n^{(a,b+1)}(z) = \frac{2(n+b+1)P_n^{(a,b)}(z)
  + 2(n+1)P_{n+1}^{(a,b)}(z)}{(2n+a+b+2)(1+z)}.
\end{equation}
Indeed, taking $z=1-2t^2$ in~\eqref{eq:Jacobirecu} and using the
definition of the generalized Gegenbauer polynomials, we have
\begin{equation}
\label{eq:gegenrec}
  (\alpha+\beta+1) (1-r^2)
  C_{n-1}^{(\beta+3/2,\alpha+1/2)} (r) = A_n
  C_{n-1}^{(\beta+1/2,\alpha+1/2)} (r) - B_n
  C_{n+1}^{(\beta+1/2,\alpha+1/2)} (r),
\end{equation}
with
\begin{align*}
A_n &=
\begin{cases}
  \dfrac{(\beta+k+1)(\alpha+\beta+k+1)}{\alpha+\beta+2k+2}, & \text{if $n=2k+1$,} \\[6pt]
  \dfrac{(\beta+k)(\alpha+\beta+k+1)}{\alpha+\beta+2k+1}, & \text{if $n=2k$,}
\end{cases} \\
B_n &=
\begin{cases}
  \dfrac{(k+1)(\alpha+k+1)}{\alpha+\beta+2k+2}, & \text{if $n=2k+1$,} \\[6pt]
  \dfrac{k(\alpha+k+1)}{\alpha+\beta+2k+1}, & \text{if $n=2k$.}
\end{cases}
\end{align*}
From~\eqref{eq:gegenrec} we have the decomposition
\begin{multline*}
  (1-r^2)K_{\beta,\alpha}(t,r)
  = \frac{1}{2(\alpha+\beta+1)^2} \sum_{n=1}^{\infty} A_n
  \frac{C_n^{(\beta+1/2,\alpha+1/2)}(t)
  C_{n-1}^{(\beta+1/2,\alpha+1/2)}(r)} {h_{n-1}^{(\beta+1,\alpha)}}
  \\
  - \frac{1}{2(\alpha+\beta+1)^2} \sum_{n=1}^{\infty} B_n
  \frac{C_n^{(\beta+1/2,\alpha+1/2)}(t)
  C_{n+1}^{(\beta+1/2,\alpha+1/2)}(r)} {h_{n-1}^{(\beta+1,\alpha)}}.
\end{multline*}
With the previous identity, we get
\begin{multline*}
  T_{\beta,\alpha}g_\lambda(t)
  = \frac1{2(\alpha+\beta+1)^2}
  \sum_{n=2}^\infty a_{n-1}(\lambda)A_n
  \frac{h_{n-1}^{(\beta,\alpha)}}{h_{n-1}^{(\beta+1,\alpha)}}
  C_n^{(\beta+1/2,\alpha+1/2)}(t)
  \\
  - \frac1{2(\alpha+\beta+1)^2}
  \sum_{n=1}^\infty a_{n+1}(\lambda) B_n
  \frac{h_{n+1}^{(\beta,\alpha)}}{h_{n-1}^{(\beta+1,\alpha)}}
  C_n^{(\beta+1/2,\alpha+1/2)}(t).
\end{multline*}
In this way, identifying the coefficients in both sides of~\eqref{eq:problem-auto},
we obtain the recurrence relation
\begin{equation*}
  \lambda a_n(\lambda)
  = \frac{1}{2(\alpha+\beta+1)^2}
  \left( a_{n-1}(\lambda) A_n
  \frac{h_{n-1}^{(\beta,\alpha)}}{h_{n-1}^{(\beta+1,\alpha)}}
  - a_{n+1}(\lambda) B_n
  \frac{h_{n+1}^{(\beta,\alpha)}}{h_{n-1}^{(\beta+1,\alpha)}} \right),
\qquad n>1,
\end{equation*}
and
\begin{equation*}
  \lambda a_1(\lambda) = -a_2(\lambda)
  \frac{B_1}{2(\alpha+\beta+1)^2}
  \frac{h_{2}^{(\beta,\alpha)}}{h_{0}^{(\beta+1,\alpha)}},
\end{equation*}
which, applying~\eqref{eq:hC2n} and~\eqref{eq:hC2n1}, becomes
\begin{equation}
\label{eq:recurrencea1}
  \lambda a_n(\lambda)
  = \frac{a_{n-1}(\lambda)}{2(\alpha+\beta+n)}
  -\frac{a_{n+1}(\lambda)}{2(\alpha+\beta+n+2)}, \qquad n>1,
\end{equation}
and
\begin{equation}
\label{eq:recurrencea2}
  \lambda a_1(\lambda) = -\frac{a_2(\lambda)}{2(\alpha+\beta+3)}.
\end{equation}

Now, we can prove the following:

\begin{Thm}
For $\alpha,\beta >-1$, and $\alpha+\beta>-1$, let
$R_{\beta+1/2,\alpha+1/2}$ be the closure of the span of
$\{C_n^{(\beta+1/2,\alpha+1/2)}\}_{n\ge 1}$ in
$L^2((-1,1),d\mu_{\beta+1,\alpha})$; then
\begin{equation}
\label{eq:decomL2}
  L^2((-1,1),d\mu_{\beta+1,\alpha})
  = R_{\beta+1/2,\alpha+1/2} \oplus R_{\beta+1/2,\alpha+1/2}^\bot,
\end{equation}
where
\begin{equation}
\label{eq:orto}
  R_{\beta+1/2,\alpha+1/2}^\bot =
\begin{cases}
  \spa\{(1-t^2)^{-1}\}, &\text{for }\, \beta>0,\\
  \{0\}, & \text{for }\, 0\ge \beta>-1.
\end{cases}
\end{equation}
Furthermore, if we let $g_\lambda(x)\in R_{\beta+1/2,\alpha+1/2}$
have the orthogonal expansion~\eqref{eq:expaneigen}, then the
eigenvalue equation~\eqref{eq:problem-auto} holds if and only if
\begin{equation}
\label{eq:conditionproblem}
  \sum_{n=1}^{\infty} |a_n(\lambda)|^2 n^{2\beta-1}<\infty.
\end{equation}
\end{Thm}

\begin{proof}
It is clear that $L^2((-1,1),d\mu_{\beta,\alpha})\subset
L^2((-1,1),d\mu_{\beta+1,\alpha})$ and $T_{\beta,\alpha}$ maps
$L^2((-1,1),d\mu_{\beta+1,\alpha})$ into
$L^2((-1,1),d\mu_{\beta,\alpha})$. Moreover, $T_{\beta,\alpha}$ is
a bounded operator and its norm is controlled by a constant $M$
given by
\[
  M^2 = \frac{1}{4(\alpha+\beta+1)^2} \, \sup_{n\ge0}
  \frac{h_{n+1}^{(\beta,\alpha)}}{h_{n}^{(\beta+1,\alpha)}}.
\]
Indeed,
\begin{align*}
  \|T_{\beta,\alpha}g\|^2_{L^2((-1,1),d\mu_{\beta,\alpha})}
  &= \sum_{n=1}^{\infty}
  \frac{|g_{n-1}|^2}{4(\alpha+\beta+1)^2}h_{n}^{(\beta,\alpha)} \\
  &\le M^2 \sum_{n=1}^{\infty}
  |g_{n}|^2h_{n}^{(\beta+1,\alpha)}
  = M^2\|g\|_{L^2((-1,1),d\mu_{\beta+1,\alpha})}.
\end{align*}
Note that $M$ is finite because
\[
  \frac{h_{2k+1}^{(\beta,\alpha)}}{h_{2k}^{(\beta+1,\alpha)}}
  = \frac{(\alpha+\beta+1)^2}{(\beta+k+1)(\alpha+k+1)}
  \qquad\text{ and }\qquad
  \frac{h_{2k}^{(\beta,\alpha)}}{h_{2k-1}^{(\beta+1,\alpha)}}
  =\frac{(\alpha+\beta+1)^2}{k(\alpha+\beta+k+1)}.
\]
In this way, we can deduce that $R_{\beta+1/2,\alpha+1/2}$ is an
invariant subspace for $T_{\beta,\alpha}$ and the
decomposition~\eqref{eq:decomL2} holds. Now, each
function $f\in R_{\beta+1/2,\alpha+1/2}^\bot$ will satisfy the
conditions $\|f\|_{L^2((-1,1),d\mu_{\beta+1,\alpha})}<\infty$ and
\begin{equation}
\label{eq:condiorto}
  \int_{-1}^{1} f(r) C_{n}^{(\beta+1/2,\alpha+1/2)}(r) 
  \, d\mu_{\beta+1,\alpha}(r) = 0,
  \qquad n=1,2,\dots.
\end{equation}
Using that $\{C_{n}^{(\beta+1/2,\alpha+1/2)}\}_{n\ge 0}$ is a
complete orthogonal system in
$L^2((-1,1),d\mu_{\beta,\alpha})$ and the relation
$d\mu_{\beta+1,\alpha}(r) = (1-r^2)\,d\mu_{\beta,\alpha}(r)$, 
from~\eqref{eq:condiorto} we get
\[
  (1-r^2)f(r) = K
\]
for a certain constant~$K$. So, taking into account that $f\in
L^2((-1,1),d\mu_{\beta+1,\alpha})$, we conclude~\eqref{eq:orto}.

Let $g_\lambda$ be a function having the expansion~\eqref{eq:expaneigen} 
and verifying~\eqref{eq:problem-auto}. Then,
using that $g_{\lambda}\in L^2((-1,1),d\mu_{\beta,\alpha})$,
it follows that
\[
  \|g_\lambda\|^2_{L^2((-1,1),d\mu_{\beta,\alpha})}
  = \sum_{n=1}^\infty |a_n(\lambda)|^2h_{n}^{(\beta,\alpha)} < \infty
\]
and this imply~\eqref{eq:conditionproblem} because
$h_{n}^{(\beta,\alpha)} \approx n^{2\beta-1}$.

On other hand, let us suppose that~\eqref{eq:conditionproblem} 
and~\eqref{eq:expaneigen} hold. To prove that $g_\lambda$ is a
solution of~\eqref{eq:problem-auto}, we need rewrite the function
$g_{\lambda}$ in terms of $C_n^{(\beta+3/2,\alpha+1/2)}$. To this
end, we use that the Jacobi polynomials satisfy the identity
\cite[formula 22.7.19, p.~782]{AbSt}
\[
  (2n+a+b)P_n^{(a,b-1)}(z) = (n+a+b)P_n^{(a,b)}(z) + (n+a)P_{n-1}^{(a,b)}(z)
\]
to produce
\[
  C_n^{(\beta+1/2,\alpha+1/2)}(t)
  = \frac{\alpha+\beta+1}{\alpha+\beta+n+1}
  (C_n^{(\beta+3/2,\alpha+1/2)}(t) - C_{n-2}^{(\beta+3/2,\alpha+1/2)}(t)).
\]
So, we find
\begin{multline}
\label{eq:glam}
  g_{\lambda}(t) = \sum_{n=1}^\infty
  \frac{\alpha+\beta+1}{\alpha+\beta+n+1}a_{n}(\lambda)
  C_n^{(\beta+3/2,\alpha+1/2)}(t) \\
  - \sum_{n=1}^\infty
  \frac{\alpha+\beta+1}{\alpha+\beta+n+1}a_{n}(\lambda)
  C_{n-2}^{(\beta+3/2,\alpha+1/2)}(t).
\end{multline}
Thus,
\[
  \|g_\lambda\|^2_{L^2((-1,1),d\mu_{\beta+1,\alpha})}
  \sim \sum_{n=1}^\infty |a_{n}(\lambda)|^2 n^{2\beta-1} < \infty
\]
and $g_\lambda \in L^2((-1,1),d\mu_{\beta+1,\alpha})$. Moreover,
from~\eqref{eq:glam}, with~\eqref{eq:recurrencea1} and~\eqref{eq:recurrencea2}, 
we can check that $g_\lambda$ is an eigenfunction of~$T_{\beta,\alpha}$.
\end{proof}

It is clear, from~\eqref{eq:recurrencea1} and~\eqref{eq:recurrencea2}, 
that $a_1(\lambda)\not=0$ (in other case $a_n(\lambda)=0$ for $n>1$). 
So this is a multiplicative factor and can be factored out. To verify the 
condition~\eqref{eq:conditionproblem}, we have to renormalize the sequence
$\{a_n(\lambda)\}_{n\ge 1}$. Set
\begin{equation}
\label{eq:relationab}
  a_n(\lambda) = i^{n-1}
  \frac{(\alpha+\beta+n+1)}{(\alpha+\beta+2)}b_{n-1}(i\lambda)a_1(\lambda).
\end{equation}
Then, using the relations~\eqref{eq:recurrencea1} and~\eqref{eq:recurrencea2}, 
we can check that
\begin{equation}
\label{eq:recurrenceb}
  2w (\alpha+\beta+n+2)b_n(w) = b_{n-1}(w)+b_{n+1}(w),
  \qquad n\ge 1,
\end{equation}
where $w=i\lambda$, $b_{-1}(w)=0$, and $b_0(w)=1$. If $R_{n,a}(z)$
denotes the Lommel polynomials, we define
$h_{n,a}(z) = R_{n,a}(1/z)$, which are known as modified Lommel
polynomials. Lommel polynomials satisfy the three terms recurrence
relation (see \cite[\S\,9.63, formula~(2), p.~299]{Wat})
\[
  \frac{2(n+a)}{z}R_{n,a}(z) = R_{n-1,a}(z)+R_{n+1,a}(z).
\]
In this way, we deduce that $b_n(w) = h_{n,\alpha+\beta+2}(w)$.

Let us identify the values verifying~\eqref{eq:conditionproblem}.

\begin{Thm}
\label{thm:eigenvalues}
The convergence condition~\eqref{eq:conditionproblem} holds 
if and only if $\lambda$ is purely imaginary, $\lambda\not=0$, and
$J_{\alpha+\beta+1}(i/\lambda) = 0$.
\end{Thm}

\begin{proof}
By Hurwitz' theorem \cite[\S\,9.65, formula~(1), p.~302]{Wat}, we have
the asymptotic relation
\[
  h_{n,a}(z) \sim \Gamma(n+a) (2z)^{n+a-1}J_{a-1}(1/z),
  \qquad n\to\infty.
\]
This fact, taking into account~\eqref{eq:relationab} and that
$b_{n}(i\lambda) = h_{n-1,\alpha+\beta+2}(i\lambda)$, shows that in
order for~\eqref{eq:conditionproblem} to hold it is necessary that
$J_{\alpha+\beta+1}(i/\lambda) = 0$ or possibly $\lambda=0$. For
$\lambda=0$ we can deduce, by~\eqref{eq:recurrenceb}, that
$b_{2n+1}(0) = 0$ and $b_{2n}(0) = (-1)^n$, then~\eqref{eq:conditionproblem} 
does not hold.

To deduce the sufficiency of $J_{\alpha+\beta+1}(i/\lambda) = 0$ we
need the identity \cite[\S\,9.6, formula~(1), p.~295]{Wat}
\[
  J_{a+n}(z) = R_{n,a}(z)J_{a}(z)-R_{n-1,a+1}(z)J_{a-1}(z).
\]
From this relation we have
\begin{equation}
\label{eq:JH}
  -J_{\alpha+\beta+n+1}(i/\lambda)
  = h_{n-1,\alpha+\beta+2}(\lambda/i) J_{\alpha+\beta}(i/\lambda)
\end{equation}
when $J_{\alpha+\beta+1}(i/\lambda) = 0$. In this way, by the
asymptotic
\[
  J_{a+n}(z) \sim \frac{(2z)^{-n-a}}{\Gamma(a+n+1)},
  \qquad n\to\infty,
\]
we can conclude that~\eqref{eq:conditionproblem} is satisfied.
\end{proof}

Finally, we can obtain the eigenvalues and the eigenfunctions of~$T_{\beta,\alpha}$
(remember that the Bessel function $J_a$ has an increasing
sequence of positive zeros, and that we are denoting them by $j_{a,k}$, for $k\ge 1$).

\begin{Thm}
\label{thm:main}
Let $\alpha,\beta >-1$, and $\alpha+\beta>-1$. Then the
eigenvalues of the integral operator $T_{\beta,\alpha}$ are $\{\pm
i/j_{\alpha+\beta+1,k}\}_{k\ge 1}$ and the eigenfunctions $g_{\pm
i/j_{\alpha+\beta+1,k}}(t)$ have the series expansion in terms of
the generalized Gegenbauer polynomials
\begin{equation}
\label{eq:expanGegen}
  \sum_{n=1}^\infty (\mp i)^{n-1}
  \frac{(\alpha+\beta+n+1)}{(\alpha+\beta+2)}
  h_{n-1,\alpha+\beta+2} \left(\frac{1}{j_{\alpha+\beta+1,k}}\right)
  C_n^{(\beta+1/2,\alpha+1/2)}(t).
\end{equation}
Moreover,
\begin{equation}
\label{eq:eigenfin}
  g_{\pm i/j_{\alpha+\beta+1,k}}(t)
  = \mp i \left(\frac{j_{\alpha+\beta+1,k}}{2}\right)^{\alpha+\beta+1}
  \frac{E_\alpha(\mp t j_{\alpha+\beta+1,k})}{\Gamma(\alpha+\beta+1)(\alpha+\beta+2)
  J_{\alpha+\beta}(j_{\alpha+\beta+1,k})}.
\end{equation}
\end{Thm}

\begin{proof}
The eigenvalues follow immediately from Theorem~\ref{thm:eigenvalues}.
The expression~\eqref{eq:expanGegen} for the eigenfunctions is a 
consequence of~\eqref{eq:expaneigen}, \eqref{eq:relationab}, the fact
$b_n(i\lambda) = h_{n-1,\alpha+\beta+2}(i\lambda)$ and the identity
\begin{equation*}
  h_{n-1,\alpha+\beta+2} \left(i\frac{\pm i}{j_{\alpha+\beta+1,k}}\right)
  = (\mp 1)^{n-1}h_{n-1,\alpha+\beta+2} \left(\frac{1}{ j_{\alpha+\beta+1,k}}\right),
\end{equation*}
which is obtained using that $R_{n,a}(-z)=(-1)^nR_{n,a}(a)$.

Let us prove~\eqref{eq:eigenfin}. Taking $\lambda=\pm
i/j_{\alpha+\beta+1,k}$ in~\eqref{eq:JH}, we obtain that
\[
  J_{\alpha+\beta+n+1}(j_{\alpha+\beta+1,k})
  = - h_{n-1,\alpha+\beta+2} \left(\frac{1}{j_{\alpha+\beta+1,k}}\right)
  J_{\alpha+\beta}(j_{\alpha+\beta+1,k}).
\]
From the previous identity, \eqref{eq:expanGegen} becomes
\begin{multline*}
  g_{\pm i/j_{\alpha+\beta+1,k}}(t)
  = \frac{-1}{(\alpha+\beta+2)J_{\alpha+\beta}(j_{\alpha+\beta+1,k})}\\*
\begin{aligned}
  &\kern20pt \times\sum_{n=1}^\infty (\mp i)^{n-1} (\alpha+\beta+n+1)
  J_{\alpha+\beta+n+1}(j_{\alpha+\beta+1,k})
  C_n^{(\beta+1/2,\alpha+1/2)}(t) \\
  &= \frac{\mp i\,j_{\alpha+\beta+1,k}^{\alpha+\beta+1}}
  {(\alpha+\beta+2)J_{\alpha+\beta}(j_{\alpha+\beta+1,k})} \\*
  &\kern20pt \times \sum_{n=0}^\infty i^{n}(\alpha+\beta+n+1)
  \frac{J_{\alpha+\beta+n+1}(\mp j_{\alpha+\beta+1,k})}
  {(\mp j_{\alpha+\beta+1,k})^{\alpha+\beta+1}}
  C_n^{(\beta+1/2,\alpha+1/2)}(t) \\
  &= \mp i \left( \frac{j_{\alpha+\beta+1,k}}{2}\right)^{\alpha+\beta+1}
  \frac{E_\alpha(\mp t j_{\alpha+\beta+1,k})}
  {\Gamma(\alpha+\beta+1)(\alpha+\beta+2) J_{\alpha+\beta}(j_{\alpha+\beta+1,k})}
\end{aligned}
\end{multline*}
where in the last step we have used~\eqref{eq:JA-Dunkl}.
\end{proof}

%------------------------------
\begin{thebibliography}{99}
%------------------------------

\bibitem{AbSt} M. Abramowitz and I. A. Stegun (Editors),
% \textit{Handbook of mathematical functions},
\textit{Handbook of mathematical functions with formulas, graphs
and mathematical tables}, United States Government Printing
Office, 1964.
Reprinted by Dover, 1972 (10th reprint). Also available on
\url{http://www.math.sfu.ca/~cbm/aands/}

\bibitem{ACV} L. D. Abreu, \'O. Ciaurri, and J. L. Varona,
Bilinear biorthogonal expansions and the Dunkl kernel on the real line,
\textit{Expo. Math.} \textbf{30} (2012), 32--48.

\bibitem{AJ} N. B. Andersen and M. de Jeu, {Elementary proofs of
Paley-Wiener theorems for the Dunkl transform on the real line},
\textit{Int. Math. Res. Not.} \textbf{2005}, no.~30, 1817--1831.

\bibitem{BCV} J. Betancor, \'O. Ciaurri, and J. L. Varona, {The multiplier
of the interval $[-1,1]$ for the Dunkl transform on the real line,}
\textit{J. Funct. Anal.} \textbf{242} (2007), 327--336.

\bibitem{IJacob} J. S. Christiansen and M. E. H. Ismail,
{A moment problem and a family of integral evaluations},
\textit{Trans. Amer. Math. Soc.}  \textbf{358} (2006), 4071--4097.

\bibitem{CV} \'O. Ciaurri and J. L. Varona, {A Whittaker-Shannon-Kotel'nikov
sampling theorem related to the Dunkl transform,} \textit{Proc. Amer. Math.
Soc.} \textbf{135} (2007), 2939--2947.

\bibitem{Dunkl} C. F. Dunkl, {Differential-difference operators associated
with reflections groups,} \textit{Trans. Amer. Math. Soc.} \textbf{311}
(1989), 167--183.

\bibitem{Dunkl2} C. F. Dunkl, {Integral kernels with reflections group
invariance,} \textit{Canad. J. Math.} \textbf{43} (1991), 1213--1227.

\bibitem{DX} C. F. Dunkl and Y. Xu,
\textit{Orthogonal polynomials of several variables},
Encyclopedia of Mathematics and its Applications No.~81,
Cambridge University Press, 2001.

\bibitem{Ism} M. E. H. Ismail, \textit{Classical and
quantum orthogonal polynomials in one variable},
Encyclopedia of Mathematics and its Applications No.~98,
Cambridge University Press, 2005.

\bibitem{IR} M. E. H. Ismail and M. Rahman,
{An inverse to the Askey-Wilson operator},
\textit{Rocky Mountain J. Math.} \textbf{32} (2002), 657--678.

\bibitem{IRZ} M. E. H. Ismail, M. Rahman, and R. Zhang,
{Diagonalization of certain integral operators~II},
\textit{J. Comp. Appl. Math.} \textbf{68} (1996), 163--196.

\bibitem{ISem} M. E. H. Ismail and P. C. Simeonov,
The spectrum of an integral operator in weigthed $L_2$ spaces,
\textit{Pacific J. Math.} \textbf{198} (2001), 443--476.

\bibitem{IZ} M. E. H. Ismail and R. Zhang, {Diagonalization of certain
integral operators}, \textit{Adv. Math.} \textbf{109} (1994), no.~1, 1--33.

\bibitem{Jeu} M. F. E. de Jeu, The Dunkl transform, \textit{Invent. Math.}
\textbf{113} (1993), 147--162.

\bibitem{Rosenblum} M. Rosenblum, {Generalized Hermite polynomials and the
Bose-like oscillator calculus,} \textit{Oper. Theory Adv. Appl.} \textbf{73}
(1994), 369--396.

\bibitem{Rosler} M. R\"osler, {A positive radial product formula
for the Dunkl kernel,} \textit{Trans. Amer. Math. Soc.} \textbf{355}
(2003), 2413--2438.

\bibitem{Wat} G. N. Watson, \textit{A Treatise on the Theory of
Bessel Functions} (2nd edition), Cambridge Univ. Press,
Cambridge, 1944.

\end{thebibliography}


%------------------------------
\end{document}
%------------------------------
